\section*{Chapitre \ref{ch:g12:space} --- Vecteurs, droites et plans dans l'espace}

\begin{solution}{exo:g12:space:1}
$\vect{AB}(2, 1, -1)$ et $\vect{AC}(1, -1, 1)$ ne sont pas proportionnels
($\frac21 \neq \frac{1}{-1}$), donc les points ne sont pas alignés. Droite
$(AB)$~:
\[
(x, y, z) = (1 + 2t,\; t,\; 2 - t), \qquad t \in \R .
\]
\end{solution}

\begin{solution}{exo:g12:space:2}
$2(x - 1) - (y - 1) + 3(z - 0) = 0$, \emph{c.-à-d.}
\[
2x - y + 3z - 1 = 0 .
\]
Pour $B(0,2,1)$~: $0 - 2 + 3 - 1 = 0$~: oui, $B$ est sur le plan.
\end{solution}

\begin{solution}{exo:g12:space:3}
$\vect{AB}(1,1,0)$, $\vect{AC}(1,0,1)$~:
\[
\cos\widehat{A}
= \frac{\vect{AB}\cdot\vect{AC}}{\norm{\vect{AB}}\,\norm{\vect{AC}}}
= \frac{1}{\sqrt2\,\sqrt2} = \frac12,
\]
donc $\widehat A = 60^\circ$.
\end{solution}

\begin{solution}{exo:g12:space:4}
Points de $\mathcal D$~: $(1 + t,\, 2 - t,\, 2t)$. En substituant dans
l'\omterm{def:g10:algebra:equation}{équation} de $\mathcal P$~:
\[
2(1+t) + (2 - t) - 2t + 1 = 5 - t = 0 \iff t = 5 .
\]
Point d'\omterm{def:g10:numbers:interunion}{intersection} unique~: $(6, -3, 10)$.
\end{solution}

\begin{solution}{exo:g12:space:5}
Les \omterm{def:g12:space:normal}{vecteurs normaux} $(1, -1, 2)$ et $(2, 1, 1)$ ne sont pas \omterm{def:g11:vect:collinear}{colinéaires},
donc les plans s'intersectent selon une droite. En additionnant les deux
\omterm{def:g10:algebra:equation}{équations}~: $3x + 3z = 5$, donc $x = \frac53 - z$~; puis la première
\omterm{def:g10:algebra:equation}{équation} donne $y = x + 2z - 1 = \frac23 + z$. Avec $z = t$~:
\[
(x, y, z) = \left(\frac53 - t,\; \frac23 + t,\; t\right), \qquad t \in \R,
\]
une droite de \omterm{def:g11:vect:direction}{vecteur directeur} $(-1, 1, 1)$. (Les deux \omterm{def:g10:algebra:equation}{équations} se
vérifient identiquement.)
\end{solution}

\begin{solution}{exo:g12:space:6}
Les directions $\vec u_1(1, 2, -1)$ et $\vec u_2(1, 1, 2)$ ne sont pas
\omterm{def:g11:vect:collinear}{colinéaires}, donc les droites ne sont pas parallèles. Une \omterm{def:g10:numbers:interunion}{intersection}
exigerait
\[
1 + t = 2 + s, \qquad 2t = 1 + s, \qquad 3 - t = 1 + 2s .
\]
Les deux premières donnent $t = 1 + s$ et $2(1+s) = 1 + s$, donc
$s = -1$, $t = 0$~; mais alors la troisième lit $3 = -1$, absurde. Pas de
point commun et non parallèles~: les droites sont non \omterm{def:g12:space:collinear}{coplanaires}.
\end{solution}

\begin{solution}{exo:g12:space:7}
\emph{1.} $\operatorname{dist} = \dfrac{\abs{2\cdot3 - 1 + 2\cdot1 - 3}}
{\sqrt{4 + 1 + 4}} = \dfrac{4}{3}$.

\emph{2.} $H = M_0 + t\,\vec n$ avec $\vec n(2, -1, 2)$, en choisissant
$t$ pour que $H \in \mathcal P$~:
\[
2(3 + 2t) - (1 - t) + 2(1 + 2t) - 3 = 4 + 9t = 0 \iff t = -\frac49 .
\]
D'où $H\left(\frac{19}{9}, \frac{13}{9}, \frac19\right)$, et
\[
M_0H = \norm{t\,\vec n} = \frac49 \times 3 = \frac43,
\]
en accord avec la formule de distance.
\end{solution}

\begin{solution}{exo:g12:space:8}
\omterm{def:g10:coordgeom:system}{Coordonnées}~: $G(1,1,1)$, et le plan $(BDE)$ passe par $B(1,0,0)$,
$D(0,1,0)$, $E(0,0,1)$.

\emph{1.} Le plan $(BDE)$ a pour \omterm{def:g10:algebra:equation}{équation} $x + y + z = 1$ (satisfaite par
les trois points, et ils ne sont pas alignés), donc $\vec n(1,1,1)$ lui
est normal. Comme $\vect{AG} = (1,1,1) = \vec n$, la diagonale $(AG)$ est
\omterm{def:g11:scal:orthogonal}{orthogonale} au plan $(BDE)$.

\emph{2.} Distance de $A(0,0,0)$~:
$\frac{\abs{0 + 0 + 0 - 1}}{\sqrt3} = \frac{1}{\sqrt3} = \frac{\sqrt3}{3}$.
La droite $(AG)$ est $(t, t, t)$~; elle rencontre le plan lorsque
$3t = 1$~: au point $\left(\frac13, \frac13, \frac13\right)$, le
centre de gravité du triangle $BDE$.
\end{solution}

\begin{solution}{exo:g12:space:9}
\emph{1.} $\vect{AB}(1, 0, -1)$, $\vect{AC}(-1, 1, 0)$,
$\vect{AD}(1, 1, 1)$. Si $\vect{AD} = a\vect{AB} + b\vect{AC}$, les
\omterm{def:g10:coordgeom:system}{coordonnées} donnent $a - b = 1$, $b = 1$, $-a = 1$~: les deux premières
forcent $a = 2$, contredisant $a = -1$. Non \omterm{def:g12:space:collinear}{coplanaires}.

\emph{2.} $\vect{BC}(-2, 1, 1)$, $\vect{BD}(0, 1, 2)$. Résoudre
$\vec n \cdot \vect{BC} = \vec n \cdot \vect{BD} = 0$~: de
$b + 2c = 0$, $b = -2c$~; puis $-2a - 2c + c = 0$ donne $c = -2a$. Avec
$a = 1$~: $\vec n(1, 4, -2)$. Plan passant par $B(2,1,0)$~:
\[
x + 4y - 2z - 6 = 0 .
\]

\emph{3.} Distance de $A(1,1,1)$~:
$\dfrac{\abs{1 + 4 - 2 - 6}}{\sqrt{21}} = \dfrac{3}{\sqrt{21}}$.
Aire de $BCD$~: $\norm{\vect{BC}}^2 = 6$, $\norm{\vect{BD}}^2 = 5$,
$\vect{BC}\cdot\vect{BD} = 3$, donc
\[
\text{Aire} = \frac12\sqrt{6 \times 5 - 9} = \frac{\sqrt{21}}{2}.
\]

\emph{4.}
\[
V = \frac13 \times \frac{\sqrt{21}}{2} \times \frac{3}{\sqrt{21}}
= \frac12 .
\]
\end{solution}

\begin{solution}{pb:g12:space:1}
\textbf{1.} Plan~: $x + y + z = 1$. Droite~: $(t, t, t)$.
\omterm{def:g10:numbers:interunion}{Intersection}~: $3t = 1$, d'où
$\left(\frac13, \frac13, \frac13\right)$.

\textbf{2.} $\dfrac{\abs{0 + 0 + 0 - 1}}{\sqrt3} =
\dfrac{1}{\sqrt3} = \dfrac{\sqrt3}{3}$.

\textbf{3.} $1 - 1 + 0 = 0$~: ils sont \omterm{def:g11:scal:orthogonal}{orthogonaux}.

\textbf{4.} Même \omterm{def:g12:space:normal}{vecteur normal} $(1,1,1)$, mais
$\frac52 \neq 1$~: plans strictement parallèles. Pour la droite~:
$(1,1,0) \cdot (1,1,1) = 2 \neq 0$, donc elle n'est pas parallèle
au plan et le perce en un point.

\textbf{5.} Milieux~: $\left(1, \frac12, 0\right)$,
$\left(\frac12, 1, 0\right)$, $\left(0, 1, \frac12\right)$,
$\left(0, \frac12, 1\right)$, $\left(\frac12, 0, 1\right)$,
$\left(1, 0, \frac12\right)$~: chacun a pour somme de \omterm{def:g10:coordgeom:system}{coordonnées}
$\frac32$.

\textbf{6.} Le \omterm{def:g12:space:normal}{vecteur normal} de $P$ est $(1,1,1)$, la direction de
$(AG)$~: ils sont perpendiculaires. Le centre
$\left(\frac12,\frac12,\frac12\right)$ a pour somme de \omterm{def:g10:coordgeom:system}{coordonnées}
$\frac32$~: il appartient à $P$.

\textbf{7.} Deux milieux consécutifs diffèrent de \omterm{def:g11:vect:vector}{vecteurs} du type
$\left(-\frac12, \frac12, 0\right)$~: de longueur
$\frac{\sqrt2}{2}$ à chaque fois --- six côtés égaux.

\textbf{8.} En $\left(\frac12, 1, 0\right)$~: \omterm{def:g11:vect:vector}{vecteurs} d'arêtes
$\left(\frac12, -\frac12, 0\right)$ et
$\left(-\frac12, 0, \frac12\right)$~: \omterm{def:g12:space:dot}{produit scalaire}
$-\frac14$, \omterm{def:g11:scal:dot}{normes} $\frac{\sqrt2}{2}$, donc
$\cos = -\frac12$ et l'angle vaut $120^\circ$. Côtés égaux, angles
égaux à $120^\circ$~: c'est un hexagone régulier.

\textbf{9.} Périmètre $6 \times \frac{\sqrt2}{2} = 3\sqrt2
\approx 4.24$. Aire $6 \times \frac{\sqrt3}{4} \times
\frac12 = \frac{3\sqrt3}{4} \approx 1.30$~: plus grande qu'une face
(d'aire $1$) --- une tranche plus large que n'importe quel côté de
la boîte~; seul le rectangle diagonal ($\sqrt2 \approx 1.41$) fait
mieux.

\textbf{10.} La droite $(t,t,t)$ rencontre $P$ en $t = \frac12$~:
c'est le centre --- lequel est bien le \omterm{prop:g10:coordgeom:midpoint}{milieu} de $[AG]$ ($A$ en
$t = 0$, $G$ en $t = 1$).

\textbf{11.} Pour $c = \frac12$~: le plan coupe les trois arêtes
issues du coin $A$ en $\left(\frac12,0,0\right)$,
$\left(0,\frac12,0\right)$ et $\left(0,0,\frac12\right)$~: un
triangle équilatéral de côté $\frac{\sqrt2}{2}$. Quand $c$ croît,
le triangle grandit, ses coins se tronquent en un hexagone
(régulier exactement en $c = \frac32$), puis la figure se rétracte
symétriquement jusqu'à un triangle au coin $G$~: la métamorphose du
tailleur de cristal.

\textbf{12.} Chaque côté de la section est l'\omterm{def:g10:numbers:interunion}{intersection} du plan de
coupe avec une \emph{face} du cube, et un plan rencontre une face
(un polygone plat) en au plus un segment~: il y a donc au plus $6$
côtés. Sept est impossible~; trois, quatre, cinq et six se
produisent tous (les questions 8 et 11 en montrent deux).

\textbf{13.} Base $ABD$~: triangle rectangle d'aire $\frac12$~;
hauteur $AE = 1$~: volume
$\frac13 \times \frac12 \times 1 = \frac16$.

\textbf{14.} Deux points quelconques parmi $B, D, E, G$ diffèrent
d'exactement $1$ sur deux \omterm{def:g10:coordgeom:system}{coordonnées}~: les six distances valent
toutes $\sqrt2$ --- le tétraèdre est régulier. Le cube se
décompose en $BDEG$ plus quatre tétraèdres de coin isométriques à
$ABDE$~: volume $1 - 4 \times \frac16 = \frac13$.

\textbf{15.} Plan $(BDE)$~: $x + y + z = 1$~; centre
$\left(\frac12,\frac12,\frac12\right)$~: distance
$\frac{\abs{\frac32 - 1}}{\sqrt3} = \frac{1}{2\sqrt3} =
\frac{\sqrt3}{6}$.

\textbf{16.} $\vect{AG} = (1,1,1)$ et $\vect{BH} = (-1,1,1)$~:
$\cos\theta = \frac{-1 + 1 + 1}{3} = \frac13$, donc
$\theta \approx 70.5^\circ$, de supplémentaire environ
$109.5^\circ$. Quatre doublets électroniques se repoussent et
s'écartent le plus possible autour du carbone~: l'optimum est
constitué des sommets du tétraèdre régulier, dont les directions
issues du centre se rencontrent exactement sous cet angle --- le
diamant est la question 14 cristallisée.

\textbf{17.} $(1, t, t)$ appartient à $P$ lorsque
$1 + 2t = \frac32$, soit $t = \frac14$~: le point
$\left(1, \frac14, \frac14\right)$.

\textbf{18.} $(AB)$~: les points $(t, 0, 0)$~; $(EG)$~: les points
$(s, s, 1)$. Se rencontrer exigerait $0 = 1$ sur la troisième
coordonnée~: jamais. Être parallèles exigerait $(1,0,0)$ et
$(1,1,0)$ \omterm{def:g11:vect:collinear}{colinéaires}~: non plus. Elles ne se coupent pas et
ne sont pas parallèles~: elles sont non
\omterm{def:g12:space:collinear}{coplanaires} --- deux couloirs à des étages différents.

\textbf{19.} $B + t(1,1,1) = (1 + t, t, t)$ appartient à $P$
lorsque $1 + 3t = \frac32$, soit $t = \frac16$~: le projeté est
$\left(\frac76, \frac16, \frac16\right)$, dont la somme des
\omterm{def:g10:coordgeom:system}{coordonnées} vaut $\frac32$~: il est bien sur $P$.

\textbf{20.} Les droites paramétrées percent (questions 1, 10 et
17)~; les \omterm{def:g12:space:normal}{vecteurs normaux} mesurent angles et distances (questions
2, 8, 15 et 16)~; les \omterm{def:g10:algebra:equation}{équations} de plans taillent les sections (la
métamorphose du triangle à l'hexagone). Et le terrain de jeu a bien
rendu la visite~: un hexagone régulier dans une boîte de carrés, un
tétraèdre régulier dans les coins, l'angle du diamant, et deux
droites qui s'ignorent --- de la géométrie que seul l'espace peut
offrir.
\end{solution}
